% !TEX program = lualatex
% !TEX root = main.tex
\documentclass[a4paper,11pt]{ltjsarticle}

% ---------- レイアウト・タイポグラフィ ----------
\usepackage[margin=25mm]{geometry}
\usepackage{microtype}

% ---------- 数式 ----------
\usepackage{amsmath, amssymb, mathtools}
\usepackage{amsthm}

% ---------- 図表 ----------
\usepackage{graphicx}
\usepackage{booktabs}
\usepackage{caption}
\usepackage{subcaption}
\usepackage{xcolor}

% ---------- リスト・引用・コード ----------
\usepackage{enumitem}
\usepackage{csquotes}
\usepackage{listings}

% ---------- 相互参照・リンク（最後に読み込む） ----------
\usepackage{hyperref}
\hypersetup{
  colorlinks=true,
  linkcolor=blue!60!black,
  urlcolor=blue!60!black,
  citecolor=blue!60!black,
}

% ---------- 定理環境 ----------
\theoremstyle{plain}
\newtheorem{theorem}{定理}[section]
\newtheorem{lemma}[theorem]{補題}
\newtheorem{proposition}[theorem]{命題}
\theoremstyle{definition}
\newtheorem{definition}[theorem]{定義}
\newtheorem{example}[theorem]{例}
\theoremstyle{remark}
\newtheorem*{remark}{注意}
\renewcommand{\proofname}{証明}

% ---------- リスティングのスタイル ----------
\lstdefinestyle{code}{
  basicstyle=\ttfamily\small,
  keywordstyle=\color{blue!60!black}\bfseries,
  commentstyle=\color{gray},
  stringstyle=\color{teal},
  numbers=left,
  numberstyle=\tiny\color{gray},
  showstringspaces=false,
  breaklines=true,
  frame=single,
  framerule=0.5pt,
  rulecolor=\color{gray!50},
}
\lstset{style=code}

% ---------- 独自ショートカット ----------
\newcommand{\R}{\mathbb{R}}
\newcommand{\N}{\mathbb{N}}
\newcommand{\abs}[1]{\left\lvert #1 \right\rvert}
\newcommand{\norm}[1]{\left\lVert #1 \right\rVert}

\title{タイトル}
\author{著者名\\\small \texttt{your.email@example.com}}
\date{\today}

\begin{document}
\maketitle

\begin{abstract}
ここに概要を書きます。問題設定・アプローチ・主要な結果を 2〜3 文でまとめます。
\end{abstract}

\tableofcontents
\bigskip

\section{はじめに}
\label{sec:intro}

段落は空行で区切ります。\emph{強調}や\textbf{太字}も使えます。
脚注\footnote{脚注は同じページの下部に表示されます。}も自然に扱えます。

外部リンクはそのまま貼れます: \url{https://example.org}。
別表記のリンクも可能です: \href{https://example.org}{サンプルサイト}。

\section{数式}
\label{sec:math}

文中の数式は $a^2 + b^2 = c^2$ のように地の文に溶け込みます。
別行立ては独立した行に表示されます:
\begin{equation}
  \label{eq:euler}
  e^{i\pi} + 1 = 0.
\end{equation}
式~\eqref{eq:euler}のように参照できます。

整列した式:
\begin{align}
  (x + y)^2 &= x^2 + 2xy + y^2, \\
  (x - y)^2 &= x^2 - 2xy + y^2.
\end{align}

場合分け:
\begin{equation}
  \abs{x} =
  \begin{cases}
    x  & (x \ge 0 \text{ のとき}), \\
    -x & (\text{それ以外}).
  \end{cases}
\end{equation}

行列:
\begin{equation}
  A =
  \begin{pmatrix}
    a_{11} & a_{12} \\
    a_{21} & a_{22}
  \end{pmatrix},
  \qquad
  \det A = a_{11} a_{22} - a_{12} a_{21}.
\end{equation}

\subsection{定理と証明}

\begin{definition}
  関数 $f \colon \R \to \R$ が点 $x_0 \in \R$ で\emph{連続}であるとは、
  任意の $\varepsilon > 0$ に対してある $\delta > 0$ が存在して、
  $\abs{x - x_0} < \delta$ ならば $\abs{f(x) - f(x_0)} < \varepsilon$ が成り立つことをいう。
\end{definition}

\begin{theorem}[ピタゴラスの定理]
  \label{thm:pythagoras}
  直角三角形の 2 辺の長さを $a$, $b$、斜辺を $c$ とすると、
  \[
    a^2 + b^2 = c^2.
  \]
\end{theorem}

\begin{proof}
  証明は省略する。標準的な参考文献~\cite{euclid}を参照のこと。
\end{proof}

\begin{remark}
  定理~\ref{thm:pythagoras}は内積空間における高次元化が知られている。
\end{remark}

\section{リスト・引用・コード}

\subsection{リスト}

箇条書き:
\begin{itemize}
  \item 1 つめの項目。
  \item 2 つめの項目（入れ子にもできる）:
  \begin{itemize}
    \item ネストした項目。
  \end{itemize}
\end{itemize}

ラベルを変えた番号付きリスト:
\begin{enumerate}[label=(\alph*)]
  \item 手順 1。
  \item 手順 2。
\end{enumerate}

定義リスト:
\begin{description}
  \item[定義域] 入力の集合。
  \item[値域] 出力になりうる値の集合。
\end{description}

\subsection{ブロック引用}

\begin{displayquote}
  数学とは、神が宇宙を記述するために用いた言語である。
  \hfill --- ガリレオ・ガリレイ（伝）
\end{displayquote}

\subsection{コードリスティング}

\begin{lstlisting}[language=Python, caption={短い例。}, label={lst:example}]
def fib(n):
    """n 番目のフィボナッチ数を返す。"""
    a, b = 0, 1
    for _ in range(n):
        a, b = b, a + b
    return a
\end{lstlisting}

リスト~\ref{lst:example}は反復版のヘルパーを示す。

\section{図と表}

\subsection{図}

\begin{figure}[ht]
  \centering
  % \includegraphics[width=0.6\linewidth]{example}
  \fbox{\rule{0pt}{6em}\rule{0.6\linewidth}{0pt}}
  \caption{プレースホルダーは \texttt{\textbackslash includegraphics} で置き換えてください。}
  \label{fig:example}
\end{figure}

\subsection{2 枚並べた図}

\begin{figure}[ht]
  \centering
  \begin{subfigure}[t]{0.45\linewidth}
    \centering
    \fbox{\rule{0pt}{4em}\rule{0.9\linewidth}{0pt}}
    \caption{左の図。}
    \label{fig:left}
  \end{subfigure}\hfill
  \begin{subfigure}[t]{0.45\linewidth}
    \centering
    \fbox{\rule{0pt}{4em}\rule{0.9\linewidth}{0pt}}
    \caption{右の図。}
    \label{fig:right}
  \end{subfigure}
  \caption{2 枚の小図からなる図。}
  \label{fig:subs}
\end{figure}

図~\ref{fig:example}と、図~\ref{fig:subs}の各パネル
(\subref{fig:left}, \subref{fig:right}) を参照のこと。

\subsection{表}

\begin{table}[ht]
  \centering
  \caption{booktabs スタイルの表。}
  \label{tab:example}
  \begin{tabular}{lrr}
    \toprule
    項目     & 件数 & 平均 (秒) \\
    \midrule
    手法 A   &  120 & 0.42 \\
    手法 B   &   85 & 0.31 \\
    手法 C   &  210 & 0.58 \\
    \midrule
    合計     &  415 & 0.44 \\
    \bottomrule
  \end{tabular}
\end{table}

表~\ref{tab:example}は 3 手法の比較である。

\section{まとめ}

主な貢献を要約し、今後の課題を述べる。

\section*{謝辞}

LaTeX エコシステムを支える OSS コミュニティに感謝する。

\begin{thebibliography}{9}
  \bibitem{euclid}
    Euclid, \emph{Elements}, ca.\ 300 BCE.
  \bibitem{knuth1984}
    D.~E.\ Knuth, \emph{The \TeX book}, Addison--Wesley, 1984.
\end{thebibliography}

\end{document}
